\section{Angular Power Spectrum}
A scalar quantity (it is ok for now to consider scalar quantities, such as matter overdensity $\delta(x)$), can be expressed in terms of spherical harmonics \cite{Chiarenza_et_al_2024}. \\
A scalar quantity $\tilde{A}(\hat{n})$ with $\hat{n}$ describes the unit vector which is in direction of the sky can be related to a 3D quantity $A(\textbf{r},z)$ through a projection kernel $W_i(z)$:
\begin{equation}
    \tilde{A}(\hat{n}) = \int dz W_i(z) A(\chi(z) \hat{n}; z)
\label{eq1_blast}
\end{equation}
where:
\begin{itemize}
    \item $\tilde{A}(\hat{n})$ is the quantity we observe in the sky, and is a function of the unit vector $\hat{n}$ (\textbf{\textit{Domanda: it is a 2D vector, is it not?}}). It seems to me that this is a 2D field on the sphere: for every direction $\hat{n}$, we have a single number; 
    \item $\int dz $ is the integral on all the redshifts - that is, from $z = 0$ up to the depth our survey reaches; 
    \item $W_i(z)$ is called \textit{projection kernel}, and weighs how much of the 3D field contributes to the observation, at every redshift. For example, for galaxy clustering, $W_i(z)$ has a peak near the redshift of our galaxy sample, because we are interested only in galaxies at certain redshifts; instead for weak lensing the $W_i(z)$ peaks at intermediate redshifts, because we are interested in structures between distant galaxies and our local universe;
    \item $A(\textbf{r},z)$ is our 3D cosmological field; it can be for example the density contrast $\delta_m(r)$. It is defined at physical coordinates $\vec{r} = \chi(z)\hat{n}$, and at redshift $z$;
    \item $\chi$ is the comoving distance;
\end{itemize}
This quantity can be decomposed in its "constituent frequencies": following the definition of Fourier transform \ref{fourier_transforms}, one can do the same for the quantity $A({\textbf{r}}; z) = A((\chi(z)\hat{\textbf{n}}),z)$:
\begin{equation}
    A(\chi(z) \hat{\textbf{n}}; z) = \int_{-\infty}^{\infty} \frac{d^3k}{(2\pi)^3} A(\vec{\textbf{k}}, z) e^{i\vec{\textbf{k}}\cdot\hat{\textbf{n}}\chi(z)} 
\label{eq2_blast}
\end{equation}
where we decompose the field $A(\chi(z) \hat{\textbf{n}}; z)$ in simple plane waves\footnote{where 
\begin{itemize}
    \item $A(\chi(z) \hat{\textbf{n}}; z)$ is the field in the real space at position $\chi(z) \hat{\textbf{n}}$; 
    \item $\frac{1}{(2\pi)^3}$ is a normalization factor coming from Parseval's theorem; 
    \item $A(\vec{\textbf{k}}; z)$ is the amplitude of Fourier's modes: is a complex number, and represents the amplitude and the phase of every plane wave. Here $\vec{\textbf{k}}$ is a 3D wavevector, and represents the wavelength and the oscillation direction. Here small $|{\vec{\textbf{k}}}|$ describe large scales, and large $|{\vec{\textbf{k}}}|$ describe small scales;
    \item $e^{i\vec{\textbf{k}}\cdot\hat{\textbf{n}}\chi(z)}$ is the plane wave with wavevector ${\vec{\textbf{k}}}$;
\end{itemize}
}, where every plane wave has a wavevector $\vec{\textbf{k}}$ and contains a specific amplitude, $A(\vec{\textbf{k}}; z)$:
The idea will be to further decompose the plan waves in spherical harmonics, in order to isolate angular scales on the sky. \\
The plane wave can be decomposed then in spherical harmonics converting the cartesian coordinates in spherical coordinates, using Rayleigh expansion \footnote{
\begin{equation}
\label{eq3_blast}
    e^{i\vec{\textbf{k}}\cdot\hat{\textbf{n}}r} = 4\pi \sum_{l,m} i^l j_l(kr) Y^{*}_{lm}(\hat{\textbf{k}}) Y_{lm}(\hat{\textbf{n}})
\end{equation}
where:
\begin{itemize}
    \item $e^{i\vec{\textbf{k}}\cdot\hat{\textbf{n}}r}$ is the plane wave with direction $\hat{\textbf{n}}$ and distance $r = \chi(z)$;
    \item 4$\pi$ is a normalization factor;
    \item $\sum_{l,m}$ is the sum on all the modes of the multipole. Here $l = 0, 1, ..., \infty$ and $m = -l, ... 0, ..., l$. The order of multipole $l$ is related to the angular scales on the sky by  \textit{resolution} $\approx 180°/l$, while the $m$ describes the azimuthal modes, that encode the rotational dependency around a symmetry axis;
    \item $i^l$ is a phase convention;
    \item $j_{lm}$ is the spherical Bessel function; is the oscillation radial profile for every angular mode $l$. It oscillates with frequency $k$, and peaks near $r \approx l/k$. This means that modes with greater $l$ are concentrated to bigger radii;
    \item $Y^{*}_{lm}(\hat{\textbf{k}}) Y_{lm}(\hat{\textbf{n}})$ are the spherical harmonics; $Y^{*}_{lm}(\hat{\textbf{k}})$ projects the direction $\hat{\textbf{k}}$ on the multipole $(l,m)$, while $Y_{lm}(\hat{\textbf{n}})$ does the same for direction $\hat{\textbf{n}}$, They are analogue to sines and cosines;
\end{itemize}
} one separates the plane wave in independent angular modes, each with a different radial profile. \\
This means that angular scales (different $l$) are decoupled: this allows the computation of $C_l$ separately for each $l$. A plane wave with high $k$ contains the contributes of all $l$: by measuring different modes $l$, one can sample the same structure at different angular scales (\textbf{\textit{Domanda: is it like shooting the same subject with different lenses, i.e. a fisheye, a portrait mode, a landscape, a wide-angle lens?}}). \\ 
Combining \ref{eq2_blast} and \ref{eq3_blast} one gets:
\begin{equation}
    A(\chi(z) \hat{\textbf{n}}; z) = 4\pi \sum_{l,m} i^l \int \frac{d^3k}{(2\pi)^3} A(\vec{\textbf{k}}, z) j_l(k\chi(z)) Y^{*}_{lm}(\hat{\textbf{k}}) Y_{lm}(\hat{\textbf{n}}) 
\label{eq4_blast}
\end{equation}
and combining \ref{eq1_blast} with \ref{eq4_blast} one gets:
\begin{equation}
    \tilde{A}(\hat{n}) = 4\pi \sum_{l,m} i^l \int dz W_i(z) \int \frac{d^3k}{(2\pi)^3} A(\vec{\textbf{k}}, z) j_l(k\chi(z)) Y^{*}_{lm}(\hat{\textbf{k}}) Y_{lm}(\hat{\textbf{n}})
\label{eq5_blast}
\end{equation}
with:
\begin{equation}
    A_{l,m} = 4\pi i^l \sum_{l,m} \int dz W_i(z) \int \frac{d^3k}{(2\pi)^3} A(\vec{\textbf{k}}, z) j_l(k\chi(z)) Y^{*}_{lm}(\hat{\textbf{k}})
\label{eq6_blast}
\end{equation}
coefficients of the spherical harmonics decomposition\footnote{
where:
\begin{itemize}
    \item $A_{l,m}$ is a single complex number representing the amplitude of the multipole mode ($l,m$). One extracts the amplitude of the spherical harmonics from the 3D field;
    \item $\int dz W_i(z)$ weighs every redshift with the kernel function, that selects which "slice" of the Universe contributes;
    \item $\int \frac{d^3k}{(2\pi)^3} A(\vec{\textbf{k}}, z)$ is the integration on all the Fourier's modes $k$;
    \item $j_l(k\chi(z))$ connects Fourier's wavenumber $k$ to the multipole $l$ through the Bessel function, at distance $\chi(z)$. The Bessel function automatically encode the link between the Fourier scales and the angular scales;
    \item $Y^{*}_{lm}(\hat{\textbf{k}})$ projects the direction of Fourier's wavevector $\vec(\textbf{k})$ on the selected multipole ($l,m$). The spherical harmonics select which angular scales are of interest;
    \end{itemize}
}. \\
This is how one describes a single scalar field $\tilde{A}(\hat{\textbf{n}})$. Defining a second scalar field $\tilde{B}(\hat{\textbf{n}})$, one can compute the \textbf{cross-correlation power spectrum}\footnote{
where:
\begin{itemize}
    \item $C_{AB}(l)$ is the covariance of the two fields. That is, the variance of the $l$-th angular mode in the cross-correlation of fields $A$ and $B$;
    \item $\frac{2}{\pi}$ is a normalization that comes when combining $\frac{(4\pi)(4\pi)(2\pi)^3}{(2\pi)^3)(2\pi)^3)}$;
    \item $\int dz W_i(z) \int dz' W_j(z')$ are the integrations on the redshifts for the two fields. $dz$ and $dz'$ link different "times", while $W_i(z)$ and $W_j(z')$ project the 3D fields in the tomographic bins of redshifts $z$ and $z'$;
    \item $ \int dk$ is the integration on the wavenumbers;
    \item $P_{AB}(k,z,z')$ is the correlation of the Fourier modes at different redshifts $z, z'$. If $A=B=cluster$, this represents the squared matter power spectrum; if the fields represent the cluster and shear probes, then it represents the cross-correlation of those different physical quantities;
    \item $j_l(k\chi(z)) j_l(k\chi(z'))$ is a pair of spherical Bessel functions that couple the Fourier scale $k$ to the angular scale $l$, for the probes at redshifts $z$ and $z'$. The product of the two Bessel functions is non-zero when both $k\chi(z)$ and $k\chi(z')$ are near $(l+\frac{1}{2})$, that is the peak of $j_l$. This determines which wavenumbers and distances contribute the most at each $l$. These functions operate the angular projection from 3D to 2D;
    \end{itemize}
}:
\begin{equation}
    C_{AB}(l) = \langle A_{l,m} B_{l,m}^*\rangle = \frac{2}{\pi} \int dz W_i(z) \int dz' W_j(z') \int dk k^2 P_{AB}(k,z,z') j_l(k\chi(z)) j_l(k\chi(z'))
\label{eq7_blast}
\end{equation}
The window functions $W_i(z)$ and $W_j(z')$ depend on the specific probe - that is, in our case, cluster or shear. \\
Moreover, one can operate a change of variables for the integral in $z$ and $z'$, to $\chi$ and $\chi'$. This leads to the final form of the angular power spectrum\footnote{
where:
\begin{itemize}
    \item $C_{AB}(l)$ is the angular power spectrum with tomographic bins $i,j$;
    \item $N(l)$ is a normalization which depends on $l$ and on the probe used;
    \item $\int d\chi_1 \int d\chi' \int dk$ is the triple integration, the reason to have an approximation, and the culprit of it all;
    \item $W_i^A(\chi_1)$ and $W_j^B(\chi_2)$ are the new kernel functions. $W_i^A(\chi_1)$ represents the weight for the field $A$, in the $i$-th bin, at distance $\chi_1$. The shape of the kernels depends on the probe:
    \begin{itemize}
        \item for clustering: \begin{equation}
                                W_i^g(\chi) = H(z)n_i(z)b_g(z)
                                \label{kernel_cl}
                              \end{equation}
        where $H(z)$ is the Hubble factor, $n_i(z)$ is the galaxy number density, and $b_g(z)$ is the galaxy bias;
        \item for weak-lensing: \begin{equation}
                                W_i^s(\chi) = \frac{3H_0^2\Omega_m}{2a}\chi\int_{z}^{\infty}dz'n_i(z')\frac{\chi(z')-\chi(z)}{\chi(z')}
                                \label{kernel_sh}
                              \end{equation}
    \end{itemize}
    \item $P^{AB}(k, \chi_1, \chi_2)$ is the 3D matter power spectrum;
    \item $j_l(k\chi_1)$ and $j_l(k\chi_2)$ are the two spherical Bessel functions;
    \item $(k\chi_1)^\alpha$ and $(k\chi_2)^\beta$ are the prefactors depending on the probes: for clustering, there are no radial dependencies other that the Bessel functions, and $\alpha = \beta = 0$, whereas for shear $\alpha = \beta = 2$ - this is due to the weak-lensing convergence which is a projected quantity;
    \end{itemize}
}:
\begin{equation}
    C_{ij}^{AB}(l) = N(l) \int d\chi_1 W_i^A(\chi_1) \int d\chi_2 W_j^B(\chi_2) \int dk k^2 P_{AB}(k,\chi_1,\chi_2) \frac{j_l(k\chi_1) j_l(k\chi_2)}{(k\chi_1)^\alpha(k\chi_2)^\beta}
\label{eq8_blast}
\end{equation}
which is the challenging integral to compute. \\
Let's write the three angular power spectra for the possible combination of probes. \\
For galaxy clustering ($gg$), the power spectra is:
\begin{equation}
    C_{ij}^{gg}(l) = \frac{2}{\pi} \int_0^{\infty} d\chi_1 W_i^g(\chi_1) \int_0^{\infty} d\chi_2 W_j^g(\chi_2) \int_0^{\infty} dk k^2 P_{AB}(k,\chi_1,\chi_2) j_l(k\chi_1) j_l(k\chi_2)
\label{eq15_blast}    
\end{equation}
where $P^{gg}(k,\chi_1, \chi_2)$ is the matter power spectrum and describes the amplitudes of density fluctuations at wavenumber $k \approx 1/l$. \\ 
The matter power spectrum $P^{gg}(k,\chi)$ relates to the density fluctuations through the linear galaxy bias $b_g(\chi)$, viz: $P^{gg}(k,\chi) = b_g(\chi)^2 P_{\delta}(k, \chi)$. \\
For weak lensing ($ss$), the power spectra is:
\begin{equation}
        C_{ij}^{ss}(l) = \frac{2(l+2)!}{\pi(l-2)!} \int_0^{\infty} d\chi_1 W_i^s(\chi_1) \int_0^{\infty} d\chi_2 W_j^s(\chi_2) \int_0^{\infty} dk k^2 P_{AB}(k,\chi_1,\chi_2) \frac{j_l(k\chi_1) j_l(k\chi_2)}{(k\chi_1)^2(k\chi_2)^2}
\label{eq17_blast}    
\end{equation}
where in \ref{kernel_sh} the fraction outside the integral is related to the "strength" of the lens, $\chi$ is the comoving distance of the probe (far-away structures lens more), $P^{ss}(k,\chi_1, \chi_2)$ does not contain bias, like the galaxy power spectrum. \\ 
The term $(k\chi)^2$ at the denominator in the integral comes from the spin-like nature of lensing: it's a second derivative of the lensing potential which returns $(ik)(ik) = -k^2$. \\
It represents the fact that modes at large scales (small $l$) produce small shear, whereas modes at small scales (large $l$) produce large shear. \\
For the cross-correlation ($gs$), the power spectra is:
\begin{equation}
        C_{ij}^{gs}(l) = \frac{2}{\pi} \sqrt\frac{(l+2)!}{(l-2)!} \int_0^{\infty} d\chi_1 W_i^g(\chi_1) \int_0^{\infty} d\chi_2 W_j^s(\chi_2) \int_0^{\infty} dk k^2 P_{AB}(k,\chi_1,\chi_2) \frac{j_l(k\chi_1)}{(k\chi_1)^2}j_l(k\chi_2)
\label{eq19_blast}    
\end{equation}
The thing is: the oscillating integral is the inner one, the one in $k$. \\
The other two, those in $\chi_1$ and $\chi_2$, are solvable with simple (and fast) quadrature rules. 
The approach of the paper is to approximate the power spectrum through a decomposition with Chebyshev polynomials, obtaining\footnote{
where: 
    \begin{itemize}
        \item $c_n(\chi_1, \chi_2; \theta)$ are the coefficients depending on the distance pairs $\chi_1$ and $\chi_2$, and $\theta$, which contain the cosmology parameters. These coefficients are computed only once;
        \item $T_n(k)$ are the Chebyshev polynomials.
    \end{itemize}
The approximation used is the following: a function $f(x)$ can be approximated using interpolation at specific points, Chebyshev points, defined by: 
\begin{equation}
    t_n = cos(\frac{2n+1}{2N}\pi) 
\label{eq25_blast}
\end{equation}
with $n = 0,1,...,N-1$. These points are defined in $[-1,1]$, and are denser near the edge of the interval. The interpolating polynomial $\hat{f}(x) \approx f(x)$ is:
\begin{equation}
    \hat{f}(x) = \sum_{n=0}^{N-1} c_n T_n(x)
\label{eq26_blast}
\end{equation}
where $c_n$ can be obtained via a FFT transform:
\begin{equation}
    c_n = \sum_{j=0}^{N-1}f(x_j)e^{2i\pi j n/N}
\label{eq29_blast}
\end{equation}
}:
\begin{equation}
    P(k,\chi_1, \chi_2;\theta) \approx \sum_{n = 0}^{n_{max}} c_n(\chi_1, \chi_2; \theta) T_n(k)
\label{eq30_blast}
\end{equation}
The integral in $dk$ in \ref{eq8_blast} gets the approximation:
\begin{equation}
    w_l^{AB}(\chi_1,\chi_2;\theta) \equiv \sum_{n = 0}^{n_{max}} c_n(\chi_1, \chi_2; \theta) \tilde{T}_{n;l}^{AB}(\chi_1,\chi_2)
\label{eq31_blast}
\end{equation}
where\footnote{
where:
\begin{equation}
        f^{AB}(k) = \begin{cases} 
                k^2 & AB = gg \\
                1/k^2 & AB = ss \\
                1 & AB = gs 
                \end{cases}
    \label{eq33_blast}
    \end{equation}
and $w_l^{AB}$ is the projected matter density. The $\tilde{T}_{n;l}^{AB}(\chi_1,\chi_2)$ is computed once for every pair $(n,l)$ and then stored. The cosmology is contained in the $\theta$ of the $c_n$, whereas the geometry is contained in $\tilde{T}_{n;l}^{AB}$.
}
\begin{equation}
    \tilde{T}_{n;l}^{AB}(\chi_1,\chi_2) \equiv \int_{k_{min}}^{k_{max}} dk f^{AB}(k) T_n(k) j_l(k\chi_1) j_l(k\chi_2)
\label{eq32_blast}
\end{equation}
The last thing one can do before translating in coding language the problem is to change the variables. One defines $R \equiv \chi_2/\chi_1$, obtaining a form for \ref{eq8_blast}:
\begin{equation}
    C_{ij}^{AB}(l) \approx \int d\chi_1 W_i^A(\chi_1) \int d\chi_2 {W}_j^B(\chi_2) w_l^{AB}(\chi_1,\chi_2;\theta) 
\end{equation}
that becomes after some work\footnote{
The integral of the angular power spectrum can be symmetrized. Changing the notation for the kernel from $W_i(\chi)$ to $K_i(\chi_1)$ one gets:
\begin{equation}
\begin{split}
    C_{ij}^{AB}(l) &=  \int_{0}^{\infty} d\chi_1 K_i^A(\chi_1) \int_{0}^{\infty} d\chi_2 K_j^B(\chi_1) w_{AB}(\chi_1,\chi_2) = \\
                  &\quad \int_{0}^{\infty} d\chi_1 K_i^A(\chi_1) \int_{0}^{\chi_1} d\chi_2 K_j^B(\chi_2) w_{AB}(\chi_1,\chi_2) + 
                  \int_{0}^{\infty} d\chi_2 K_j^B(\chi_2) \int_{0}^{\chi_2} d\chi_1 K_i^A(\chi_1) w_{AB}(\chi_2,\chi_1)
\end{split}
\label{eqA1_blast}
\end{equation}
in the first integral we are investigating the situation where $\chi_2 \leq \chi_1$, and we place $\chi_1 = \chi$, $R = \chi_2/\chi_1$. By computing the Jacobian, one gets that $d\chi_1 d\chi_2 = \chi d\chi dR$. In the second integral the situation is symmetric: $\chi_2 \geq \chi_1$, and we place $\chi_2 = \chi$, $R = \chi_1/\chi_2$. The Jacobian is the same. Inserting those relations into \ref{eqA1_blast}, one gets:
\begin{equation}
    C_{ij}^{AB}(l) =  \int_{0}^{\infty} d\chi \int_{0}^{1} dR \chi [\mathcal{K}_i^A(\chi) \mathcal{K}_j^B(R\chi) + \mathcal{K}_j^B(\chi) \mathcal{K}_i^A(R\chi)] 
                        w_l^{AB}(\chi,R\chi) = \int_{0}^{\infty} d\chi \int_{0}^{1} dR \chi \tilde{\mathcal{K}}_i^{AB}(\chi, R\chi)w_{AB}(\chi,R\chi)
\label{eqA2_blast}
\end{equation}
where $\mathcal{K}_i^A(\chi)$ is \ref{eq35_blast}
}:
\begin{equation}
    C_{ij}^{AB}(l) = \int_{0}^{\infty} d\chi \int_{0}^{1} dR \chi [\mathcal{K}_i^A(\chi) \mathcal{K}_j^B(R\chi) + \mathcal{K}_j^B(\chi) \mathcal{K}_i^A(R\chi)] w_l^{AB}(\chi,R\chi) 
\label{eq34_blast}
\end{equation}
with
\begin{equation}
    \mathcal{K}_i^A(\chi) = \begin{cases} 
    K_i^A(\chi) & clustering \\
    K_i^A(\chi)/\chi^2 & lensing
                \end{cases}
\label{eq35_blast}
\end{equation}